\section{Техническое задание}
\subsection{Основание для разработки}

Основанием для разработки является задание по выпускной квалификационной работе <<Поисковое приложение с графическим интерфейсом на Qt на базе библиотеки регулярных выражений>>.

\subsection{Назначение разработки}

Основными задачами выпускной квалификационной работы являются:
\begin{enumerate}
	\item Проектирование и разработка библиотеки регулярных выражений на языке Си;
	\item Проектирование и разработка графического приложения на Qt для поиска в документах с интеграцией разработанной библиотеки регулярных выражений.
	\end{enumerate}

Задачами данной разработки являются:
\begin{itemize}
	\item провести анализ предметной области;
	\item разработать концептуальную модель библиотеки регулярных выражений;
	\item спроектировать программную систему библиотеки регулярных выражений;
	\item спроектировать API библиотеки регулярных выражений для её лёгкой интеграции в программы на Си-подобных языках;
	\item реализовать библиотеку регулярных выражений средствами языка Си;
	\item спроектировать программную систему поискового приложения с графическим интерфейсом;
	\item реализовать поисковое приложение на языке C++ с использованием фреймворка Qt для графического интерфейса.

\end{itemize}

\subsection{Библиотека регулярных выражений}
\subsubsection{Функционал библиотеки регулярных выражений}

Библиотека регулярных выражений на основе строки шаблона регулярного выражения позволяет:
\begin{itemize}
	\item откомпилировать шаблон в объект регулярного выражения для дальнейшего многократного использования в обработке текста;
	\item получить вердикт, подходит ли строка под шаблон; 
	\item получить массив захваченных подстрок из заданной строки;
	\item заменить захваченные подстроки в заданной строке на указанный текст, а также получить число выполненных замен;
	\item в случае некорректности шаблона регулярного выражения получить описание допущенной ошибки.
\end{itemize}

\subsubsection{Синтаксис шаблона регулярного выражения}

Шаблон регулярного выражения должен позволять обрабатывать:
\begin{enumerate}
	\item ASCII-символы и последовательности байтов UTF-8 в буквальном виде.
	\item Метасимволы: 
	\begin{enumerate}
		\item\verb`.` – любой одиночный символ;
		\item\verb`^` – результат должен включать начало строки;
		\item\verb`$` – результат должен включать конец строки;
		\item\verb`*` – предыдущий токен (символ или группа) встречается 0 или более раз;
		\item\verb`+` – предыдущий токен встречается 1 или более раз;
		\item\verb`?` – предыдущий токен опционален;
		\item\verb`*?`, \verb`+?`, \verb`??` – ленивый режим захвата для данного токена;
		\item\verb`{n}` – предыдущий токен повторяется ровно \textit{n} раз;
		\item\verb`{n,m}` – предыдущий токен повторяется от \textit{n} до \textit{m} раз;
		\item\verb`[abc]`  – символ равен одному из символов в квадратных скобках;
		\item\verb`[^abc]`  – символ равен любому символу, которого нет в квадратных скобках;
		\item\verb`abc|def` – соответствует либо шаблону слева вертикальной черты, либо шаблону справа от нее (альтерация);
		\item\verb`(abc)` – определяет группу захвата, а также ограничивает область шаблона, рассматриваемого альтерацией;
		\item\verb`(?:abc)` – определяет не-захватываемую группу и ограничивает область шаблона, рассматриваемого альтерацией;
		\item\verb`\n` – обратная ссылка на захваченную ранее при поиске группу, где \textit{x} – номер группы от 1 до 9.
	\end{enumerate}
	\item Символьные классы:
	\begin{enumerate} 
		\item\verb`\d` – любая цифра;
		\item\verb`\D` – любой нецифровой символа;
		\item\verb`\s` – пробельный символ;
		\item\verb`\S` – любой непробельный символ;
		\item\verb`\w` – буква латинского алфавита, цифра или символ подчёркивания;
		\item\verb`\W` – любой символ, кроме буквенно-цифрового и знака подчёркивания;
		\item\verb`\a` – буква латинского алфавита;
		\item\verb`\A` – любой символ, кроме латинских букв;
		\item\verb`\p` – символы пунктуации;
		\item\verb`\P` – любой непунктуационный символ;
		\item\verb`\n` – символ новой строки;
		\item\verb`\t` – символ табуляции.
	\end{enumerate}
	\item Экранирование метасимволов: для буквальной интерпретации метасимволов, перечисленных в пункте 2, перед ними необходимо поставить экранирующий символ: $\backslash$. Сам экранирующий символ также нужно экранировать для буквальной интерпретации.
	\item Подстановка захваченной группы при замене с помощью регулярного выражения (\textbf{\$n} или \textbf{$\backslash$n}, где \textit{n} – номер группы, захваченный регулярным выражением, или всей захваченной подстроки при \textit{n} = 0).
	\item Модификаторы регулярного выражения:
	\begin{enumerate}
		\item\textit{m} – многострочный режим (\^{} и \$ означают начало и конец каждой строки);
		\item\textit{s} – метасимвол '.' включает в себя переход на новую строку;
		\item\textit{i} – поиск без учёта регистра (только латинские буквы);
		\item\textit{I} – поиск без учёта регистра (только базовая латиница и кириллица, текст обязательно должен быть в кодировке UTF-8).
	\end{enumerate}
\end{enumerate}

\subsubsection{API библиотеки регулярных выражений}

Библиотека регулярных выражений должна иметь переносимый C API для её интеграции в другие программы на Си-подобных языках посредством заголовочного (.h) файла, содержащего API-функции с пояснениями про принимаемые и возвращаемые ими параметры.

API библиотеки должно реализовывать следующие возможности:
\begin{itemize}
	\item компиляция шаблона регулярного выражения для дальнейшего повторного использования при обработке строк;
	\item валидация строки с помощью шаблона регулярного выражения;
	\item получение захваченной подстроки из заданной строки шаблона или из откомпилированного регулярного выражения;
	\item получение массива со всеми захваченными подстроками из заданной строки шаблона или из откомпилированного регулярного выражения;
	\item получение строки – результата замены захваченных подстрок (всех или до указанного ограничения на количество) в заданной строке на указанный текст, а также числа выполненных замен.
\end{itemize}

\subsubsection{Экспортируемые функции библиотеки регулярных выражений}
  % список 4-го уровня глубины, который можно использовать далее в любом месте
\newcounter{subsubsubsection}[subsubsection]
\renewcommand{\thesubsubsubsection}{\thesubsubsection.\arabic{subsubsubsection}}
\newcommand{\subsubsubsection}[1]{%
	\refstepcounter{subsubsubsection}%
	{\normalfont\normalsize\bfseries\thesubsubsubsection\quad#1}\par}
	
\subsubsubsection{Компиляция шаблона регулярного выражения}

Функции данного семейства выполняют разбор строки шаблона и создают внутреннее представление регулярного выражения в виде объекта \verb`struct re_chain`, пригодного для многократного использования без повторного разбора шаблона.

\begin{enumerate}
	\item \verb`struct re_chain* compile_re(char* pattern)` -- компиляция шаблона в объект регулярного выражения. Возвращает указатель на объект регулярного выражения; в случае некорректности шаблона возвращает NULL. Чтобы понять, в чём ошибка, нужно вызвать полную версию функции (\verb`compile_re2`) со всеми параметрами.
	
	\item \verb`struct re_chain* compile_re2(char* pattern, int* errcode, int* errpos)` -- то же, что и для функции выше, но принимает ещё 2 параметра-указателя, заполняемые в случае некорректности шаблона: \verb`errcode` -- код ошибки и \verb`errpos` -- номер символа в строке шаблона, при разборе которого произошла ошибка (может указывать на нуль-терминатор строки, если ожидался закрывающий символ). Допустимо передать NULL-указатели, тогда эти поля не будут заполнены.
	
	\item \verb`void free_re(struct re_chain* re_obj)` -- освобождение памяти, занимаемой объектом регулярного выражения, созданным функциями \verb`compile_re` и \verb`compile_re2`
\end{enumerate}

\begin{lstlisting}[caption={Пример компиляции шаблона с диагностикой ошибок}]
	char* pattern = "(\\d{2})\\.(\\d{2})\\.(\\d{4})";

	int errcode, errpos;
	struct re_chain *re = compile_re2(pattern, &errcode, &errpos);
	if (!re) {
		// errpos - позиция символа шаблона, вызвавшего ошибку
		fprintf(stderr, "Ошибка с кодом %d, позиция %d\n", errcode, errpos);
		return 1;
	}

	free_re(re); // освобождение скомпилированного объекта
\end{lstlisting}

\subsubsubsection{Поиск позиции совпадения}

Данные функции возвращают позицию и длину первого совпадения регулярного выражения в тексте, не выделяя памяти под содержимое захваченной подстроки.

\begin{enumerate}
	\item \verb`int matchpos(char* pattern, char* text, int* capture_data)` -- возвращает данные о захваченной подстроке из исходного текста по строке шаблона регулярного выражения. При успешном захвате возвращает номер первого символа подстроки и записывает её длину в переменную по указателю \verb`capture_data`, если он не равен NULL. При неудачном захвате возвращает -1. При некорректности шаблона возвращает отрицательное число меньше -1 -- код ошибки, а также записывает номер символа в строке шаблона, при разборе которого произошла ошибка, в \verb`capture_data`. Если по указателю \verb`capture_data` передать положительное число, то поиск в строке \verb`text` начнётся с этой позиции.
	
	\item \verb`int obj_matchpos(struct re_chain* re_obj, char* text, int* capture_data)` -- то же, что и для функции выше, но вместо строки с шаблоном принимает откомпилированный объект регулярного выражения.
\end{enumerate}


\begin{lstlisting}[caption={Последовательный проход по всем совпадениям в тексте}]
	char* text = "В 2023 году вышло 42 релиза, а в 2024 году - уже 58";
	struct re_chain* re = compile_re("\\d+");
	
	int pos = 0;
	int len;
	while (1) {
		len = pos;
		int start = obj_matchpos(re, text, &len); // поиск начиная с позиции len
		if (start < 0)
			break;
		
		printf("позиция %d: \"%.*s\"\n", start, len, text + start);
		pos = start + len; // сдвиг за конец найденного совпадения
	}
	free_re(re); // освобождение объекта регулярного выражения
\end{lstlisting}

\subsubsubsection{Валидация строки}

Функции данного семейства проверяют, соответствует ли строка целиком шаблону регулярного выражения, и возвращают булев результат без извлечения подстрок.

\begin{enumerate}
	\item \verb`int is_match(char* pattern, char* text)` -- выполняет валидацию строки с помощью строки шаблона регулярного выражения: при успехе возвращает 1, при неудаче -- 0, при некорректности шаблона возвращает отрицательное число меньше -1 -- код ошибки.
	
	\item \verb`int obj_is_match(struct re_chain* re_obj, char* text)` -- то же, что и для функции выше, но вместо строки с шаблоном принимает откомпилированный объект регулярного выражения.
\end{enumerate}

\begin{lstlisting}[caption={Валидация e-mail адресов}]
	struct re_chain* re_email = compile_re("^[\\w.+-]+@[\\w-]+\\.[\\w.]+$");
	char* inputs[] = {
		"user@example.com", "not-an-email",
		"wrong@abc.",     "abc@af-y.ru"
	};
	
	for (int i = 0; i < 4; i++) {
		int email_ok = obj_is_match(re_email, inputs[i]); // 1 - совпадение, 0 - нет
		printf("'%s' - %s\n", inputs[i], email_ok == 1 ? "корректен" : "не корректен!");
	}
	free_re(re_email);
\end{lstlisting}

\subsubsubsection{Извлечение одного совпадения}

Функции данного семейства возвращают первое совпадение регулярного выражения вместе с содержимым всех групп захвата в виде NULL-терминированного массива строк.

\begin{enumerate}
	\item \verb`char** match(char* pattern, char* text)` -- возвращает NULL-терминированный массив malloc-строк аналогично функции \verb`match()` в JavaScript: первый элемент -- главный захват, далее -- содержимое групп захвата. При отсутствии совпадений возвращает пустой массив, при некорректности шаблона -- NULL.
	
	\item \verb`char** obj_match(struct re_chain* re_obj, char* text)` -- то же, что и для функции выше, но вместо строки с шаблоном принимает откомпилированный объект регулярного выражения.
\end{enumerate}

\begin{lstlisting}[caption={Извлечение даты в формате ДД.ММ.ГГГГ из набора строк}]
	struct re_chain* re = compile_re("(\\d{2})\\.(\\d{2})\\.(\\d{4})");
	char* texts[] = {
		"Дата рождения: 01.04.1987",
		"Дата выдачи: 15.08.2003",
		"Дата истечения: 31.12.2030"
	};
	
	for (int i = 0; i < 3; i++) {
		char** m = obj_match(re, texts[i]); // первый захват + группы
		if (m && m[0]) {
			printf("%s: день=%s месяц=%s год=%s\n", m[0], m[1], m[2], m[3]);
			for (int j = 0; m[j]; j++) free(m[j]);
		}
		free(m);
	}
	free_re(re);
\end{lstlisting}

\subsubsubsection{Извлечение всех совпадений}

Функции данного семейства возвращают все непересекающиеся совпадения регулярного выражения в тексте в виде NULL-терминированного массива строк, аналогично вызову \verb`match()` с флагом \verb`g` в JavaScript.

\begin{enumerate}
	\item \verb`char** matches(char* pattern, char* text)` -- возвращает NULL-терминированный массив malloc-строк -- все захваченные подстроки регулярного выражения. При отсутствии совпадений возвращает пустой массив, при некорректности шаблона -- NULL.
	
	\item \verb`char** obj_matches(struct re_chain* re_obj, char* text)` -- то же, что и для функции выше, но вместо строки с шаблоном принимает откомпилированный объект регулярного выражения.
\end{enumerate}

\begin{lstlisting}[caption={Извлечение всех слов из текста}]
	// сбор всех слов из лог-строки для анализа
	char* log_line = "[ERROR] 2024-01-15 DatabaseConnection: timeout after 30s";
	struct re_chain* re = compile_re("\\a+");
	
	char** words = obj_matches(re, log_line);
	if (words) {
		printf("Найденные слова:\n");
		for (int i = 0; words[i]; i++) {
			printf("  [%d] %s\n", i, words[i]);
			free(words[i]);
		}
		free(words);
	}
	free_re(re);
\end{lstlisting}

\subsubsubsection{Замена совпадений}

Функции данного семейства выполняют замену совпадений регулярного выражения в исходной строке на строку замены, которая может содержать ссылки на группы захвата (\verb`$n` или \verb`\n`).

\begin{enumerate}
	\item \verb`int re_nreplace(char* pattern, char* repl, char* text, char* dest, int n, int limit)` -- выполняет замену подстрок, совпадающих с \verb`pattern`, на строку \verb`repl` в тексте \verb`text`, максимум \verb`limit` раз (0 -- без ограничений). Результат записывается в буфер \verb`dest` размером \verb`n` байт. Возвращает число выполненных замен, -1 при переполнении буфера или отрицательный код ошибки при некорректном шаблоне.
	
	\item \verb`int obj_re_nreplace(struct re_chain* re_obj, char* repl, char* text, char* dest,`
	  \verb`int n, int limit)` -- то же, что и для функции выше, но вместо строки с шаблоном принимает откомпилированный объект регулярного выражения.
	  
	\item \verb`char* re_dupreplace(char* pattern, char* repl, char* text, int limit, int* repl_cnt)` -- выполняет замену подстрок, совпадающих с \verb`pattern`, на строку \verb`repl` в тексте \verb`text`, максимум \verb`limit` раз (0 -- без ограничений). Результат записывается в возвращаемую malloc-строку. Записывает по указателю  \verb`repl_cnt` число выполненных замен или отрицательный код ошибки при некорректном шаблоне
	
	\item \verb`char* obj_re_dupreplace(struct re_chain* re_obj, char* repl, char* text, int limit,`
	  \verb`int* repl_cnt)` -- то же, что и для функции выше, но вместо строки с шаблоном принимает откомпилированный объект регулярного выражения.
\end{enumerate}

\begin{lstlisting}[caption={Замена подстрок с группами захвата}]
	// Приведение дат ДД.ММ.ГГГГ к формату ISO (ГГГГ-ММ-ДД)
	char* text = "Договор от 15.03.2023, срок действия до 31.12.2025";
	char dest[256];
	
	struct re_chain* re = compile_re("(\\d{2})\\.(\\d{2})\\.(\\d{4})");
	// $1 $2 $3 - ссылки на группы захвата (день, месяц, год)
	int count = obj_re_nreplace(re, "$3-$2-$1",
			text, dest, sizeof(dest), 0); // 0 - без ограничений
	if (count >= 0)
		printf("результат (%d замен(ы)): %s\n", count, dest);
	else
		fprintf(stderr, "переполнение буфера или ошибка шаблона\n");
	
	free_re(re);
	// вывод: Договор от 2023-03-15, срок действия до 2025-12-31
\end{lstlisting}

\subsection{Поисковое приложение}

Десктопное приложение для поиска и замены текста в документах должно обладать
следующими функциональными и техническими требованиями:

\begin{enumerate}
	\item Работа с файлами и текстом:
	\begin{enumerate}[label=\alph*)]
		\item позволять пользователю открывать произвольный текстовый файл
		и выполнять в нём поиск посредством регулярных выражений, подсвечивать
		все найденные совпадения и опционально выводить их в список с возможностью
		навигации по документу с его помощью;
		\item позволять выполнять замену текста в файле, используя шаблон
		регулярного выражения как в интерактивном режиме, так и в пакетном;
		\item сохранять файл после замены с возможностью перезаписать оригинал
		или сохранить результат как новый файл;
		\item отображать номера строк в редакторе, а также текущую позицию
		курсора (строка, столбец) в строке состояния;
		\item определять кодировку файла при открытии и выполнять автоматическую
		конвертацию содержимого в кодировку UTF-8 для дальнейшей работы.
	\end{enumerate}
	
	\item Поиск и навигация:
	\begin{enumerate}[label=\alph*)]
		\item иметь возможность выполнять поиск без учёта и с учётом регистра,
		выводить на экран ошибки в составлении шаблона при их наличии, а также
		выполнять простой поиск по тексту без использования регулярных выражений;
		\item отображать общее число найденных совпадений;
		\item поддерживать переход к следующему и предыдущему совпадению
		по горячим клавишам;
		\item отображать номер строки для каждого совпадения
		в списке результатов;
	\end{enumerate}
	
	\item Поиск по папке:
	\begin{enumerate}[label=\alph*)]
		\item обеспечивать возможность рекурсивного поиска по регулярному
		выражению во всех текстовых файлах в выбранной папке, выводить результаты
		поиска и при нажатии на совпадение открывать содержащий его файл
		для редактирования в поисковом приложении;
		\item поддерживать фильтрацию файлов по расширению при рекурсивном
		поиске (например, только \verb|*.txt|, \verb|*.log|);
		\item отображать число совпадений на файл в дереве результатов поиска.
	\end{enumerate}
	
	\item Пользовательский интерфейс:
	\begin{enumerate}[label=\alph*)]
		\item хранить историю последних использованных шаблонов регулярных
		выражений и предоставлять к ней доступ через выпадающий список
		поля ввода;
		\item отображать справку по синтаксису поддерживаемых регулярных
		выражений непосредственно в приложении;
		\item поддерживать светлую и темную тему оформления средствами Qt.
	\end{enumerate}
	
	\item Технические требования:
	\begin{enumerate}[label=\alph*)]
		\item быть написано с использованием фреймворка Qt на языке C++
		для обеспечения кроссплатформенности среди десктопных операционных систем;
		\item интегрировать API разработанной библиотеки регулярных выражений.
	\end{enumerate}
\end{enumerate}

\subsection{Требования к оформлению документации}

Разработка программной документации и программного изделия должна производиться согласно ГОСТ 19.102-77 и ГОСТ 34.601-90. Единая система программной документации.
